% ---------------------------------------------------------------------
% Section 5: Geometrically consistent control law and certified H-inf performance
% ---------------------------------------------------------------------
\section{Geometrically Consistent Control Law and Certified $H_\infty$ Performance}

\subsection{Disturbances: explicit solution, multiplicative/exogenous separation, assumptions}

The true joint-space dynamics with friction $\boldsymbol f$, actuation error $\delta\boldsymbol\tau$, and environmental contact $\boldsymbol\tau_{\mathrm{ext}}$ is
$M(\boldsymbol q)\ddot{\boldsymbol q}+C(\boldsymbol q,\dot{\boldsymbol q})\dot{\boldsymbol q}+\boldsymbol g(\boldsymbol q)+\boldsymbol f=\boldsymbol\tau+\delta\boldsymbol\tau+\boldsymbol\tau_{\mathrm{ext}}$;
the inner loop executes the computed torque based on the nominal model, $\boldsymbol\tau=\hat M\ddot{\boldsymbol q}_{\mathrm{ref}}+\hat C\dot{\boldsymbol q}+\hat g+\hat{\boldsymbol f}$. With the mismatch defined as ``nominal minus true'', $\Delta M\triangleq\hat M-M$ (likewise $\Delta C,\Delta\boldsymbol g,\Delta\boldsymbol f$), substituting and left-multiplying by $M^{-1}$ yields $\ddot{\boldsymbol q}=\ddot{\boldsymbol q}_{\mathrm{ref}}+\boldsymbol w_{\mathrm{dyn}}$ with
\begin{equation}
\boldsymbol w_{\mathrm{dyn}}=M^{-1}\bigl(\Delta M\,\ddot{\boldsymbol q}_{\mathrm{ref}}+\Delta C\dot{\boldsymbol q}+\Delta\boldsymbol g+\Delta\boldsymbol f+\delta\boldsymbol\tau+\boldsymbol\tau_{\mathrm{ext}}\bigr).
\tag{5.1}\label{eq:5.1}
\end{equation}

\textbf{Separation into multiplicative and exogenous parts}. Substituting $\ddot{\boldsymbol q}_{\mathrm{ref}}=J^{+}(u_{\mathrm{ff}}+u_{\mathrm{fb}}-\dot J\dot{\boldsymbol q})$ into \eqref{eq:5.1} and using $JJ^{+}=I_6$ gives the disturbance decomposition
\begin{equation}
d=\Theta(\boldsymbol q)\,u_{\mathrm{fb}}+d_{\mathrm{ex}},\qquad
\Theta\triangleq J M^{-1}\Delta M\,J^{+}\in\mathbb R^{6\times6},
\tag{5.1d}\label{eq:5.1d}
\end{equation}
where $\Theta$ is the multiplicative part and $d_{\mathrm{ex}}$ the exogenous part (see Appendix C.1). The exogenous part is further split by time profile, $d_{\mathrm{ex}}=d_{L_2}+d_b$: the former ($L_2$, noise-like) is handled by Theorem~3(c) and the latter ($L_\infty$, bias-like) by Theorem~3(d).

\textbf{Assumptions} (used silently by Theorem~3, listed here once):
\begin{itemize}[leftmargin=2em]
  \item \textbf{(A1) Regular Jacobian}: on a working set $\mathcal Q$, $J(\boldsymbol q)$ has full row rank with $\sigma_{\min}(J)\ge\underline\sigma>0$, hence $JJ^{+}=I_6$ (near singularities a damped pseudoinverse is used \cite{Nak86}, and the residual is folded into $d_{\mathrm{ex}}$).
  \item \textbf{(A2) Bounded inertia and mismatch}: $M(\boldsymbol q)\succeq \underline m I_n$ ($\underline m>0$), and $\Delta M,\Delta C,\Delta\boldsymbol g,\Delta\boldsymbol f$ are bounded on $\mathcal Q\times\{\|\dot{\boldsymbol q}\|\le\bar v\}$.
  \item \textbf{(A3) Multiplicative small-gain condition} \cite{ZDG96}:
  \begin{equation}
  \alpha\triangleq\sup_{\boldsymbol q\in\mathcal Q}\bigl\|\Theta(\boldsymbol q)\bigr\|_2
  \;<\;\frac{\lambda_{\min}(K_d)}{\lambda_{\max}(K_d)}=\frac1{\mathrm{cond}_2(K_d)}\ \le 1 .
  \tag{5.1f}\label{eq:5.1f}
  \end{equation}
  For isotropic damping this reduces to the classical robustness condition $\alpha<1$ of computed-torque control \cite{Spo92}.
  \item \textbf{(A4) Desired trajectory and exogenous signals}: $\hat{\underline x}_d(t)$ is $C^2$ with bounded $\|\boldsymbol\xi_d\|,\|\dot{\boldsymbol\xi}_d\|$; $d_{L_2}\in L_2$, $d_b\in L_\infty$, and $D_{\mathrm{ex}}\triangleq\|d_{\mathrm{ex}}\|_{L_\infty}<\infty$ (boundedness of $\dot{\boldsymbol v}_w$ requires a filtered derivative in implementation \cite{Ber93}).
\end{itemize}

\subsection{Control law}

Take $K_p$ symmetric positive definite (typically block diagonal, $K_p=\mathrm{diag}(p_O I_3,p_T I_3)$; $K_p=k_pI_6$ recovers the scalar case), and $K_d$ symmetric positive definite, block diagonal $K_d=\mathrm{diag}(K_\omega,K_v)$ when the per-channel bounds of Theorem~3(c-2) are desired. The reference acceleration (the geometrically consistent computed-torque law \cite{LWP80}) is
\begin{equation}
\ddot{\boldsymbol q}_{\mathrm{ref}}
=J^{+}\Bigl(\underbrace{\mathrm{vec}_6\bigl(\mathrm{Ad}_{\tilde x}\dot{\boldsymbol\xi}_d+\mathrm{ad}_{\tilde{\boldsymbol\xi}}\mathrm{Ad}_{\tilde x}\boldsymbol\xi_d\bigr)}_{\text{feedforward: transported desired acceleration + transport correction}}
-K_d\,e_\xi-A^{\top}(\tilde x)\,K_p\,e_z-\dot J\dot{\boldsymbol q}\Bigr).
\tag{5.2}\label{eq:5.2}
\end{equation}
The pose feedback is \emph{$A^\top$-shaped}: taking $-A^\top(\tilde x)K_pe_z$ rather than, e.g., $-K_p\log$-errors, makes the cross term in the Lyapunov derivative cancel exactly (the role of $A^\top$ is that of an energy-shaping potential expressed in the error coordinates). The feedforward and $\dot J\dot{\boldsymbol q}$ come from the $\sigma^2$ term of the desired chain and from \eqref{eq:3.5}; all feedback quantities are read off the HDQ error element of Theorem~1, and no acceleration error is measured or estimated.

\subsection{Main results}

The four results of this section (Theorems~3(a)--3(d)) share assumptions (A1)--(A4) (\S5.1), the control law \eqref{eq:5.2}, and the disturbance model \eqref{eq:5.1}/\eqref{eq:5.1d}; the common storage function is
\begin{equation}
V(e_z,e_\xi)=\tfrac12\|e_\xi\|^2+\tfrac12\,e_z^{\top}K_pe_z\;\ge0 ,
\tag{5.4a}\label{eq:5.4a}
\end{equation}
with $K_p\succ0$ symmetric ($K_p=k_pI_6$ recovers the scalar form $\tfrac{k_p}2\|e_z\|^2$). Limited by the topology of $SO(3)$ (double cover) and the singularity of $A$ at $\tilde\eta=0$, all convergence claims are local, on a sublevel set.

% ---- Theorem 3(a) ----
\setcounter{theorem}{2}%
\renewcommand{\thetheorem}{3(a)}%
\begin{theorem}[Closed-loop error dynamics]\label{thm:3a}
The error coordinates $(e_z,e_\xi)$ satisfy the cascade
\begin{equation}
\dot e_z=A(\tilde x)\,e_\xi,\qquad
\dot e_\xi=-K_d e_\xi-A^{\top}(\tilde x)\,K_p\,e_z+d(t),
\tag{5.5}\label{eq:5.5}
\end{equation}
where $d=J\boldsymbol w_{\mathrm{dyn}}+\dot{\boldsymbol v}_w+\dot{\boldsymbol v}_c$ collects all acceleration-level disturbances and decomposes exactly as $d=\Theta u_{\mathrm{fb}}+d_{\mathrm{ex}}$ by \eqref{eq:5.1d}.
\end{theorem}

\begin{proof}
Differentiating \eqref{eq:4.4} term by term and using the adjoint transport rule
$\frac{d}{dt}\bigl(\mathrm{Ad}_{\tilde x}\boldsymbol a\bigr)=\mathrm{Ad}_{\tilde x}\dot{\boldsymbol a}+\mathrm{ad}_{\tilde{\boldsymbol\xi}}(\mathrm{Ad}_{\tilde x}\boldsymbol a)$
(direct expansion from $\dot{\tilde x}=\tfrac12\tilde{\boldsymbol\xi}\tilde x$ and its conjugate) gives the twist-error dynamics
\begin{equation}
\dot{\tilde{\boldsymbol\xi}}=\dot{\boldsymbol\xi}-\mathrm{Ad}_{\tilde x}\dot{\boldsymbol\xi}_d-\mathrm{ad}_{\tilde{\boldsymbol\xi}}\bigl(\mathrm{Ad}_{\tilde x}\boldsymbol\xi_d\bigr)
\tag{5.5a}\label{eq:5.5a}
\end{equation}
(with disturbances, $\dot{\boldsymbol v}_w+\dot{\boldsymbol v}_c$ is added on the right). Substituting into this the $\dot{\boldsymbol\xi}$ obtained from \eqref{eq:5.2} via \eqref{eq:3.5}, the feedforward cancels the desired/transport terms exactly and \eqref{eq:5.5} results.
\end{proof}

% ---- Theorem 3(b) ----
\renewcommand{\thetheorem}{3(b)}%
\begin{theorem}[Nominal case: sublevel-set invariance, asymptotic and local exponential convergence]\label{thm:3b}
Let $d\equiv0$, and let $K_p=\mathrm{diag}(K_{p,O},K_{p,T})$ be symmetric positive definite (rotation/translation blocks). Take
\begin{equation}
0<c<c^{*}\triangleq\tfrac12\lambda_{\min}(K_{p,O}),
\qquad
\Omega_c\triangleq\{(e_z,e_\xi):V\le c\}.
\tag{5.5b}\label{eq:5.5b}
\end{equation}
Then: (i) $\dot V=-e_\xi^\top K_de_\xi\le0$, so $\Omega_c$ is compact and forward invariant; (ii) if $\tilde\eta(0)>0$ then along the trajectory
\begin{equation}
\tilde\eta(t)\ \ge\ \eta_0\triangleq\sqrt{1-\frac{2c}{\lambda_{\min}(K_{p,O})}}>0 ,
\tag{5.5c}\label{eq:5.5c}
\end{equation}
so $A(\tilde x)$ is uniformly invertible on $\Omega_c$ ($\det A=-\tfrac18\tilde\eta\le-\tfrac18\eta_0<0$); (iii) every trajectory in $\Omega_c$ satisfies $(e_z,e_\xi)\to(0,0)$; (iv) $(0,0)$ is \textbf{locally exponentially stable}.
\end{theorem}

\begin{proof}
The key step is the exact cancellation of the cross term $e_z^\top K_pAe_\xi$ by the $A^\top$-shaped feedback, giving $\dot V=-e_\xi^\top K_de_\xi$; LaSalle's theorem together with uniform invertibility of $A$ on $\Omega_c$ (Appendix C.2) yields asymptotic convergence. See Appendix C.6 for the full proof.
\end{proof}

\begin{remark}
The sublevel-set condition \eqref{eq:5.5b} is numerically verifiable: for the hardware tier G4 below, $c^{*}=72$ while the measured peak of $V$ is about $0.5$, a margin of more than two orders of magnitude (\S6.1). On the $\tilde\eta<0$ branch, flipping the sign of the HDQ error (which leaves $\tilde{\boldsymbol\xi}$ unchanged, Theorem~1(i)) enforces $\tilde\eta\ge0$. $\blacksquare$
\end{remark}

% ---- Theorem 3(c) ----
\renewcommand{\thetheorem}{3(c)}%
\begin{theorem}[$L_2$ disturbances: $H_\infty$ quadratic/Schur condition and per-channel split]\label{thm:3c}
Let $d=d_{L_2}\in L_2$.

\textbf{Certificate parameters}: $\kappa>0$ is the reciprocal of the error penalty (the supply-rate coefficient of $\|e_\xi\|^2$ is $\tfrac1{2\kappa}$), and $\gamma_a>0$ is the certified attenuation level; the certified quantity is the $L_2$ gain $\gamma^{*}\triangleq\gamma_a\sqrt\kappa$ from $d$ to $e_\xi$ (in seconds: $d$ is an acceleration-level and $e_\xi$ a velocity-level signal, and for zero initial state $\|e_\xi\|_{L_2}\le\gamma^{*}\|d_{L_2}\|_{L_2}$). Both are analysis parameters: they never enter the control law \eqref{eq:5.2}, but by Theorem~3(c$'$) they invert into a tuning rule---certifying a target gain $\gamma^{*}$ is equivalent to the damping lower bound $K_d\succeq(1/\gamma^{*})I_6$. In the per-channel version, $(\kappa_\omega,\gamma_\omega)$ and $(\kappa_v,\gamma_v)$ can be specified independently.

\textbf{(c-1) Pooled condition} (arbitrary symmetric positive definite $K_d$): the performance target $\dot V\le-\tfrac1{2\kappa}\|e_\xi\|^2+\tfrac{\gamma_a^2}2\|d\|^2$ holds for all $(e_\xi,d)$ \textbf{if and only if}
\begin{equation}
M\triangleq\begin{bmatrix}K_d-\tfrac1{2\kappa}I_6 & -\tfrac12 I_6\\[2pt] -\tfrac12 I_6 & \tfrac{\gamma_a^2}2 I_6\end{bmatrix}\succeq0
\;\overset{\text{Schur}}{\Longleftrightarrow}\;
K_d\succeq\tfrac12\bigl(\kappa^{-1}+\gamma_a^{-2}\bigr)I_6 ,
\tag{5.6a}\label{eq:5.6a}
\end{equation}
in which case
\begin{equation}
\int_0^\infty\kappa^{-1}\|e_\xi\|^2dt\le\gamma_a^2\int_0^\infty\|d_{L_2}\|^2dt+2V(0),
\tag{5.6}\label{eq:5.6}
\end{equation}
i.e., the $L_2$ gain from the acceleration-level disturbance to the twist-error energy is at most $\gamma_a\sqrt\kappa$ (for zero initial condition this reduces to a pure gain bound; the bound constrains $e_\xi$ only).

\textbf{(c-2) Per-channel condition}: let $K_d=\mathrm{diag}(K_\omega,K_v)$ be block diagonal and let $K_p=\mathrm{diag}(K_{p,O},\,k_{p,T}I_3)$ with the translational block isotropic (necessity in Appendix C.3). Writing $e_\xi=[\tilde\omega;\tilde v]$, $d=[d_\omega;d_v]$, and $V_\omega\triangleq\tfrac12\|\tilde\omega\|^2+\tfrac12\mathcal O^\top K_{p,O}\mathcal O$, $V_v\triangleq\tfrac12\|\tilde v\|^2+\tfrac{k_{p,T}}2\|\mathcal T\|^2$ ($V=V_\omega+V_v$), if
\begin{equation}
K_\omega\succeq\tfrac12\bigl(\kappa_\omega^{-1}+\gamma_\omega^{-2}\bigr)I_3,
\qquad
K_v\succeq\tfrac12\bigl(\kappa_v^{-1}+\gamma_v^{-2}\bigr)I_3,
\tag{5.6b}\label{eq:5.6b}
\end{equation}
then the rotation and translation channels satisfy independently
\begin{equation}
\int_0^\infty\kappa_\omega^{-1}\|\tilde\omega\|^2dt\le\gamma_\omega^2\int_0^\infty\|d_\omega\|^2dt+2V_\omega(0),
\quad
\int_0^\infty\kappa_v^{-1}\|\tilde v\|^2dt\le\gamma_v^2\int_0^\infty\|d_v\|^2dt+2V_v(0),
\tag{5.6$'$}\label{eq:5.6p}
\end{equation}
with dimensionally homogeneous quantities in each channel, and $(\kappa_\omega,\gamma_\omega)$, $(\kappa_v,\gamma_v)$ specifiable independently.
\end{theorem}

\begin{proof}
With disturbances the cross terms still cancel exactly, giving $\dot V=-e_\xi^\top K_de_\xi+e_\xi^\top d$; writing the performance target as $[e_\xi;d]^\top M[e_\xi;d]\ge0$ and taking the Schur complement yields the necessary-and-sufficient condition \eqref{eq:5.6a}; the per-channel decoupling rests on two algebraic identities ($\mathcal T\times\mathcal T=0$ and $\mathcal T\cdot(\mathcal T\times\tilde\omega)=0$). See Appendix C.7.
\end{proof}

% ---- Theorem 3(c'): optimal certificate and gamma inversion ----
\renewcommand{\thetheorem}{3(c$'$)}%
\begin{theorem}[Optimal certificate and inversion of the target gain $\gamma^{*}$]\label{thm:3c2}
Let $\gamma^{*}\triangleq\gamma_a\sqrt{\kappa}$ denote the certified $L_2$ gain and $\lambda\triangleq\lambda_{\min}(K_d)>0$. Then:
(i) (lower bound) every feasible point of the certificate family \eqref{eq:5.6a} satisfies $\gamma^{*}\ge1/\lambda$;
(ii) (attainability) $(\kappa,\gamma_a)=(\lambda^{-1},\lambda^{-1/2})$ is feasible and gives $\gamma^{*}=1/\lambda$.
Hence the optimal (smallest) gain certifiable by the family is
\begin{equation}
\gamma^{*}_{\mathrm{opt}}\;=\;\min_{(\kappa,\gamma_a)\ \mathrm{feasible}}\gamma_a\sqrt{\kappa}\;=\;\frac{1}{\lambda_{\min}(K_d)} ,
\tag{5.6c}\label{eq:5.6c}
\end{equation}
attained at $\kappa=\gamma_a^2=1/\lambda_{\min}(K_d)$, where \eqref{eq:5.6a} holds with equality. Equivalently, the \textbf{inversion rule}: a target gain $\gamma^{*}$ is certifiable \textbf{if and only if} $K_d\succeq(1/\gamma^{*})I_6$, and the leanest certificate is $(\kappa,\gamma_a)=(\gamma^{*},\sqrt{\gamma^{*}})$. The per-channel version \eqref{eq:5.6b} likewise gives $\gamma^{*}_\omega=1/\lambda_{\min}(K_\omega)$ and $\gamma^{*}_v=1/\lambda_{\min}(K_v)$.
\end{theorem}

\begin{proof}
(i)~By the AM--GM inequality (all terms positive),
$\kappa^{-1}+\gamma_a^{-2}\ \ge\ 2\sqrt{\kappa^{-1}\gamma_a^{-2}}\ =\ \dfrac{2}{\gamma_a\sqrt{\kappa}}\ =\ \dfrac{2}{\gamma^{*}}$;
substituting into \eqref{eq:5.6a} gives $\lambda\ \ge\ \tfrac12\bigl(\kappa^{-1}+\gamma_a^{-2}\bigr)\ \ge\ 1/\gamma^{*}$.
(ii)~At $\kappa=\gamma_a^2=\lambda^{-1}$ we have $\kappa^{-1}+\gamma_a^{-2}=2\lambda$, so \eqref{eq:5.6a} holds with equality and the certified gain is $\gamma_a\sqrt{\kappa}=1/\lambda$; moreover $\kappa=\lambda^{-1}>1/(2\lambda)$, so the feasible set is nonempty.
Necessity of the inversion rule is (i); sufficiency: if $\lambda\ge1/\gamma^{*}$, the point $(\kappa,\gamma_a)=(\gamma^{*},\sqrt{\gamma^{*}})$ satisfies $\kappa^{-1}+\gamma_a^{-2}=2/\gamma^{*}\le2\lambda$, is feasible, and certifies exactly $\gamma_a\sqrt{\kappa}=\gamma^{*}$.
\end{proof}

\begin{remark}
Theorem~3(c$'$) collapses the pair $(\kappa,\gamma_a)$ into a single conversion: the optimal certified gain equals $1/\lambda_{\min}(K_d)$. The tuning entry is the damping: a target attenuation level inverts into a lower bound on $K_d$; under the pole-placement parametrization of \S5.4, $\gamma^{*}=1/(a+b)$. This is equivalent to writing $\gamma$ directly into the law as in first-order $H_\infty$ schemes, except that here the knob is the physical gain $K_d$. $\blacksquare$
\end{remark}

% ---- Theorem 3(d) ----
\renewcommand{\thetheorem}{3(d)}%
\begin{theorem}[$L_\infty$ disturbances: multiplicative separation and mean-square bound on the twist error]\label{thm:3d}
Let (A1)--(A4) hold, let the disturbance decompose as in \eqref{eq:5.1d}, $d=\Theta u_{\mathrm{fb}}+d_{\mathrm{ex}}$, with $\alpha\triangleq\sup_t\|\Theta\|_2$ satisfying the small-gain condition \eqref{eq:5.1f}, and let $D_{\mathrm{ex}}\triangleq\|d_{\mathrm{ex}}\|_{L_\infty}<\infty$. Define the \textbf{effective damping}
\begin{equation}
\lambda_{\mathrm{eff}}\triangleq\lambda_{\min}(K_d)-\alpha\,\lambda_{\max}(K_d)\;>\;0
\tag{5.7a}\label{eq:5.7a}
\end{equation}
(positive by \eqref{eq:5.1f}). Suppose the trajectory remains in $\Omega_c$ of \eqref{eq:5.5b} over the interval considered (this premise only bounds $\|e_z\|$ and is verified a posteriori; see the remark below), and define the \textbf{equivalent disturbance amplitude}
\begin{equation}
D\triangleq D_{\mathrm{ex}}+\alpha\,\lambda_{\max}(K_p)\Bigl(1+\sqrt{\tfrac{2c}{\lambda_{\min}(K_{p,T})}}\Bigr)\sqrt{\tfrac{2c}{\lambda_{\min}(K_p)}} .
\tag{5.7b}\label{eq:5.7b}
\end{equation}
Then the twist error satisfies the finite-time mean-square bound
\begin{equation}
\frac1T\int_0^T\|e_\xi(t)\|^2\,dt\;\le\;\frac{D^2}{\lambda_{\mathrm{eff}}^2}+\frac{2V(0)}{\lambda_{\mathrm{eff}}\,T},
\qquad\forall T>0,
\tag{5.7c}\label{eq:5.7c}
\end{equation}
hence
\begin{equation}
\limsup_{T\to\infty}\ \mathrm{RMS}_{[0,T]}(e_\xi)
\;\triangleq\;\limsup_{T\to\infty}\Bigl(\frac1T\int_0^T\|e_\xi\|^2dt\Bigr)^{1/2}
\;\le\;\frac{D}{\lambda_{\mathrm{eff}}}
\;\xrightarrow[\ \alpha\to0\ ]{}\;\frac{\|d_{\mathrm{ex}}\|_{L_\infty}}{\lambda_{\min}(K_d)} .
\tag{5.7}\label{eq:5.7}
\end{equation}
That is, bias-type ($L_\infty$) uncertainty does not destroy boundedness; it only raises the steady mean-square level of the twist error to $D/\lambda_{\mathrm{eff}}$. The multiplicative term $\Theta u_{\mathrm{fb}}$ acts twice---it reduces the effective damping (denominator) by $\alpha\lambda_{\max}(K_d)$ and raises the equivalent disturbance amplitude (numerator) by $\alpha\lambda_{\max}(K_p)\|e_z\|$---and both effects vanish as $\alpha\to0$. The pose error $e_z$ is not covered by this theorem; its steady order of magnitude is estimated quasi-statically by the near-identity model \eqref{eq:5.9}.
\end{theorem}

\begin{proof}
Split $\Theta u_{\mathrm{fb}}$ into a damping-reduction term and an equivalent disturbance term, absorb each with Young's inequality, and integrate to obtain the RMS bound \eqref{eq:5.7}. See Appendix C.8.
\end{proof}

\begin{remark}
Since the negative term of $\dot V$ contains only $-\|e_\xi\|^2$ (the disturbance has relative degree two with respect to $e_z$), the result is an RMS bound rather than a pointwise ISS bound; strictification is required for the latter (Appendix C.4). The $\Omega_c$ premise must be verified a posteriori (\S6.1). $\blacksquare$
\end{remark}

\subsection{Near-identity linearization and the static stiffness scaling law}

Theorem~3 is qualitative and energetic; it does not return gain values. Linearizing \eqref{eq:5.5} around the origin yields a two-channel second-order model directly usable for gain tuning and experimentally falsifiable.

With $K_d=\mathrm{diag}(K_\omega,K_v)$ and $K_p=\mathrm{diag}(K_{p,O},k_{p,T}I_3)$, as $\tilde x\to1$ we have $A\to A_0=\mathrm{diag}(-\tfrac12I_3,I_3)$, and eliminating $e_\xi$ gives the two decoupled second-order equations
\begin{equation}
\ddot{\mathcal O}+K_\omega\dot{\mathcal O}+\tfrac14K_{p,O}\,\mathcal O=-\tfrac12\,d_\omega ,
\qquad
\ddot{\mathcal T}+K_v\dot{\mathcal T}+k_{p,T}\,\mathcal T=+\,d_v .
\tag{5.8}\label{eq:5.8}
\end{equation}
The rotation block $-\tfrac12I_3$ makes the effective rotational stiffness $\tfrac14K_{p,O}$; to match bandwidth across channels take $K_{p,O}=4k_{p,T}I_3$. Setting $\ddot{(\cdot)}=\dot{(\cdot)}=0$ yields the \textbf{static stiffness scaling law}
\begin{equation}
\|\mathcal T\|_{\mathrm{ss}}=\frac{\|d_v\|}{k_{p,T}} ,
\qquad
\|\mathcal O\|_{\mathrm{ss}}=\frac{2\,\|d_\omega\|}{\lambda(K_{p,O})} :
\tag{5.9}\label{eq:5.9}
\end{equation}
the steady residual is governed \emph{only} by the static stiffness, not by the damping. Pole-placement rule: for target poles $\{-a,-b\}$ per channel ($a,b>0$),
\begin{equation*}
K_\omega=K_v=(a+b)I_3,\qquad k_{p,T}=ab,\qquad K_{p,O}=4ab\,I_3 .
\end{equation*}
Combined with Theorem~3(c$'$), the optimal certified gain under this parametrization is
\begin{equation*}
\gamma^{*}=\frac{1}{a+b}:\qquad
\text{the tuning entry of }\gamma\text{ is the pole sum }a+b\ \text{(damping)},
\end{equation*}
while the steady residual is set independently by $k=ab$ (stiffness), via \eqref{eq:5.9}---each performance metric has a single dedicated knob.

Two experimentally falsifiable predictions follow from \eqref{eq:5.9}: \textbf{(P1)} multiplying the stiffness by $\rho$ divides the steady residual by $\rho$; \textbf{(P2)} the disturbance levels $(\|d_v\|,\|d_\omega\|)$ inverted from different gain tiers under the same physical conditions should coincide. \eqref{eq:5.8}--\eqref{eq:5.9} are local approximations as $\tilde x\to1$; the rigorous statements are those of Theorem~3. The verification premise of (P2)---that the task stiffness is the dominant spring---is not met under the hardware tether protocol of \S6.2, so its tether-free verification is left as future work.

Tightness (Appendix C.9): the certified value $1/\lambda_{\min}(K_d)$ equals the true $H_\infty$ gain of the closed-loop linearization, so the certificate is free of conservatism and the worst-case disturbance frequency is the channel natural frequency $\sqrt c$ ($6\,\mathrm{rad/s}$ for hardware tier G4).
